% a06_BALtorus_ct8: Step 3(F), fourth part: the field-dependent part of Ba2 Prop. 2.1, (2.57) (the collar quartic
% and the remainder O((L^k eps)^{kappa_0})|Col|); the g_k-gated terms of P^(k) of Ba2 p. 579 and the final step (2.116)
% are in the part files sec/a06_pgate_a.tex, _b.tex, _c.tex, input at the end of this file.

\subsubsection*{The collar: the field-dependent part of \cite[Prop.~2.1, (2.57)]{Ba2}}

\emph{Setting.} Step $k$ in scale-$k$ units ($t\defeq t_k$, $\eta\defeq L^{-k}$, $e_k$, $\lambda_k$, $p=p(t)$, $r=r(t)$); $\Omega_A\defeq B^k(\Lambda^{(k-1)'}_2)$, $\Omega_1\defeq B^k(\Lambda^{(k-1)'}_7)$, $\Omega'_1\defeq B^k(\Lambda^{(k)}_7)$, $\mathrm{Col}\defeq\Omega_1\setminus\Omega'_1$, $|\mathrm{Col}|=|\Lambda^{(k-1)'}_7\cap\Lambda^{(k)c}_7|$ unit cubes; $\tilde B\defeq A^{(k)}$, which equals $\theta_kA^{(k)}$ on $\Omega_A$; $\Phi\defeq\phi^{(k)}=a_kG_k(\Omega_A,\tilde B)Q_k^*(\tilde B)\Lambda^{(k-1)'}_6\varphi$ \cite[(2.56)]{Ba2}; $c_k(\Omega')$ the field-independent part of $P^{(k)}(\Omega',\tilde B,\varphi)$. Under \cite[(2.55)]{Ba2}, \cite[Prop.~2.1, (2.57)]{Ba2} asserts $P^{(k)}(\Lambda^{(k-1)'}_7,\theta_kA^{(k)},\varphi)=-\lambda_k\sum_{x\in B^k(\mathrm{Col})}\eta^d|\Phi(x)|^4+P^{(k)}(\Lambda^{(k)}_7,\theta_kA^{(k)},\varphi)+O((L^k\eps)^{\kappa_0})|\mathrm{Col}|$, ``a corollary of the analysis of the perturbation expansions'' (the sketch is \cite[p.~421]{Ba3}). Inputs: Proposition~1 with the vertex sizes of \cref{thm:a06e:T0}; \cite[(2.24)--(2.25)]{Ba1} at points of $\Omega_1$ (depth $\ge r-2\ge R_0$); \cite[(2.55), (2.60)]{Ba2} and the H\"older clause of \cite[p.~412]{Ba3} for $\tilde B$; \cite[Lemma~2.4, (2.65)--(2.66)]{Ba2} at every $y\in\Lambda^{(k-1)'}_7$ (\cref{thm:a06:BALA}(c)); the slice kernels of \cref{lem:a06e:SL} at $A=0$; the layer distances $\dist(\Lambda^{(k)}_{i+j},\Lambda^{(k)c}_i)\ge j(r-1)$, $\dist(\Omega_1,\Omega_A^c)\ge r-2$.

\begin{theorem}[the collar quartic and its remainder]\label{thm:a06f:C2}
$P^{(k)}(\Omega_1)-P^{(k)}(\Omega'_1)=-\lambda_k\sum_{x\in\mathrm{Col}}\eta^3|\Phi(x)|^4+R_{\mathrm{field}}+[c_k(\Omega_1)-c_k(\Omega'_1)]$ with $|R_{\mathrm{field}}|\le C_\beta t^{1/2-\beta}|\mathrm{Col}|$ for every $\beta>0$ and $|c_k(\Omega_1)-c_k(\Omega'_1)|\le Ct(1+\ln t^{-1})^5|\mathrm{Col}|\le C_\beta t^{1-\beta}|\mathrm{Col}|$, under the caps of Step~3(D$_2$) (Geometry). So \cite[(2.57)]{Ba2} holds with $\kappa_0=\frac12-\beta$, which implies the printed $\frac{4-d}4-\beta=\frac14-\beta$. The constants depend on $(d,L,M,\alpha_0,\bar n,a,c_1,b_0,\bar p,R,\varrho,e,\lambda,N,\beta)$, on the $O(1)$ of Proposition~1 and on $c_0,\delta_0,R_0$ of \cite[(2.24)]{Ba1}; they are free of $\eps$, $k$, $\eps_0$, the sets and $\varphi$.
\end{theorem}

The proof is \cref{lem:a06f:C21,lem:a06f:C22,lem:a06f:C23,lem:a06f:C24,lem:a06f:C25,lem:a06f:C26,lem:a06f:C27}.

\begin{lemma}[how $\Omega_1$ enters]\label{lem:a06f:C21}
In \cite[(1.4)--(1.5)]{Ba3}, $\Omega_1$ enters only through: \textup{(i)} $\mathbf 1_{\Omega_1}$ in the sums of the vertices \cite[(1.6)--(1.7)]{Ba3}; \textup{(ii)} $g_k$ ($\supp g_k\subset\Omega_1$, $g_k=1$ at distance $\ge M$ from $\partial\Omega_1$, derivatives up to order $2$ bounded by $1$) on every $A'$-leg, through $e'g_kA'+A^{(k)}$, the $A'$-legs of the pair vertices \cite[(1.14)]{Ba3} included; \textup{(iii)} $\dm_k(e'e_kg_k,\lambda'\lambda_k,\Omega_1,x)$ and $E_k(\dots,\Omega_1)$ of \cite[(1.29)--(1.30)]{Ba3}, sums of graphs whose internal vertices are localised by \textup{(i)}--\textup{(ii)} (\cite[pp.~419--420]{Ba3}: ``We localize the vertices (1.6) and (1.7) restricting summations to the set $\Omega_1$. The other vertices are localized by functions $g_k$''; ``We perform this localization for the terms of $\dm_k$, $E_k$ also''). The Gaussian measures, the propagators $G_k(\Omega_A,\tilde B)$ and the external fields $\Phi$, $\varphi$ are $\Omega_1$-free.
\end{lemma}

\begin{lemma}[marking]\label{lem:a06f:C22}
Put $D_{\mathrm{good}}\defeq\{x\in\Omega'_1:\dist(x,\partial\Omega'_1)\ge M+2\}$ and $\widetilde{\mathrm{Col}}\defeq\Omega_1\setminus D_{\mathrm{good}}$, and localise every graph in unit cubes \cite[p.~420]{Ba3}. \textup{(a)} A localised family all of whose cubes (counterterm-internal vertices included) lie in $D_{\mathrm{good}}$ has the same value for $\Omega_1$ and $\Omega'_1$. \textup{(b)} $P(\Omega_1)-P(\Omega'_1)$ is the sum over the localised families with a cube meeting $\widetilde{\mathrm{Col}}$ of $E(\cdot;\Omega_1)-E(\cdot;\Omega'_1)$. \textup{(c)} $|\widetilde{\mathrm{Col}}|\le(1+(2M+7)^d)|\mathrm{Col}|$.
\end{lemma}

\begin{proof}
$\dist(\Omega'_1,\Omega_1^c)\ge7(r-1)-1>M+3$, so $\partial\Omega'_1$ faces only $\mathrm{Col}$. On $D_{\mathrm{good}}$ (with the side-$2$ supports of the leg partitions) $\mathbf 1_{\Omega_1}=\mathbf 1_{\Omega'_1}=1$ and $g^{\Omega_1}=g^{\Omega'_1}=1$; \cref{lem:a06f:C21} for the counterterms. (c): every unit block of $\Omega'_1\setminus D_{\mathrm{good}}$ is within $M+3$ of a block of $\mathrm{Col}$, and each block of $\mathrm{Col}$ is so near to at most $(2M+7)^d$ of them.
\end{proof}

\begin{lemma}[external fields]\label{lem:a06f:C23}
For localisations $h$ at cubes in $\Omega_1$: $\norm{h\Phi}_{1,\alpha_0}\le C_\Phi p\lambda_k^{-1/4}$, and the block fields \cite[(1.12)]{Ba3} obey $|\varphi(y)|\le C_\Phi p\lambda_k^{-1/4}$ for $y\in\Lambda^{(k-1)'}_7$, with $C_\Phi\defeq3C_hac_0c_1L^{1/4}(1+\ln L)^{\bar p}$ ($C_h$ the localisation constant, $3$ for the three parts of \cite[(1.32)]{Ba3}). There are no external vector fields ($n'=0$ \cite[p.~414]{Ba3}; every $A'$-leg is contracted).
\end{lemma}

\begin{proof}
$\Phi=a_kG_k(\Omega_A,\tilde B)f$ with $f\defeq Q_k^*(\tilde B)\Lambda'_6\varphi$, $\norm f_\infty\le\sup_{\Lambda'_6}|\varphi|\le c_1p(t/L)\lambda(t/L)^{-1/4}$ by \cite[(2.55)]{Ba2} ($\Lambda'_6\subset\Lambda'_{-1}$); $p(t/L)\le(1+\ln L)^{\bar p}p(t)$, $\lambda(t/L)=\lambda_k/L$; apply \cite[(2.24)--(2.25)]{Ba1} at points of $\Omega_1$, at depth $\ge r-2\ge R_0$ in $\Omega_A$, $\tilde B$ regular by \cite[(2.60)]{Ba2}.
\end{proof}

\begin{lemma}[census and parity]\label{lem:a06f:C24}
The expressions of $P^{(k)}$ are the vertices \cite[(1.6), (1.7), (1.8), (1.10)]{Ba3} and the pairs from \cite[(1.12)--(1.14)]{Ba3}, with $n'=0$ (\cite[p.~414]{Ba3}); a pair is the scalar product of two of \cite[(1.12)--(1.14)]{Ba3}, at least one of the form \cite[(1.14)]{Ba3}, with factor $-a_k$ or $-\frac12a_k$, and \cite[(1.12)]{Ba3} is an external field; the $R$-vertices do not occur (\cite[(1.5)]{Ba3} has order $\le\bar n$). Let $s(v)$ be the number of scalar legs of $v$ (\cite[(1.12)]{Ba3} counted) and $X(v)\defeq2d_v(v)+4d_s(v)-s(v)$: $X=0$ for \cite[(1.6)]{Ba3}; $X\ge2$ for \cite[(1.7)]{Ba3} (coefficient of order $\ge2$); $X=2n-2$ for \cite[(1.8)]{Ba3} and \cite[(1.10)]{Ba3} at order $n$; $X=2m-2$ for a pair, $m\ge1$ the total order of its \cite[(1.14)]{Ba3} members. For a connected graph $G$ with $l_s$ internal scalar lines and $n_{\mathrm{ext}}$ external scalar legs, $W(G)\defeq2d_v+4d_s-n_{\mathrm{ext}}=\sum_vX(v)+2l_s$ is even and $\ge0$, and $W=0$ exactly for the lone \cite[(1.6)]{Ba3} vertex with four external legs and for two order-$1$ objects (\cite[(1.8)]{Ba3} with $n=1$, or a pair \cite[(1.12)]{Ba3}$\cdot$\cite[(1.14)]{Ba3}$_1$, \cite[(1.13)]{Ba3}$\cdot$\cite[(1.14)]{Ba3}$_1$ with both scalar legs external) joined by one vector line.
\end{lemma}

\begin{proof}
$\sum_vs(v)=n_{\mathrm{ext}}+2l_s$. $W=0$ forces $l_s=0$ and $X(v)=0$ for every $v$; the $X=0$ objects carry $0$ or $1$ $A'$-leg, all $A'$-legs are contracted and $P^{(k)}$ is connected, and a connected graph with vector lines only and at most one $A'$-leg per vertex has at most two vertices.
\end{proof}

\begin{lemma}[the families with $W\ge2$]\label{lem:a06f:C25}
For a localised family with $W(G)\ge2$, at $\Omega'\in\{\Omega_1,\Omega'_1\}$: $|E(G_{\mathrm{ren}},\{\square\};\Omega')|\le O(1)e^{d_v}\lambda^{d_s-n_{\mathrm{ext}}/4}C_\Phi^{n_{\mathrm{ext}}}p^{n_{\mathrm{ext}}}t^{W/4}e^{-\delta_0d(\{\square\})}$; summed over the marked families of \cref{lem:a06f:C22}, this part of $R_{\mathrm{field}}$ is $\le C_\beta t^{1/2-\beta}|\mathrm{Col}|$.
\end{lemma}

\begin{proof}
Proposition~1 with the external scalar fields of \cref{lem:a06f:C23} and no external vector field; the background satisfies the hypotheses of \cite[p.~412]{Ba3} ($|\partial\tilde B|\le O(p)$ by \cite[(2.60)]{Ba2}, its H\"older clause), there is no $\tilde A$, and its $O(1)$ is independent of $\eps$, $k$, the domains and $\tilde B$ \cite[p.~421]{Ba3}. In $d=3$, $e_k^{d_v}\lambda_k^{d_s}(p\lambda_k^{-1/4})^{n_{\mathrm{ext}}}=e^{d_v}\lambda^{d_s-n_{\mathrm{ext}}/4}p^{n_{\mathrm{ext}}}t^{W/4}$, and $W\ge2$ gives $t^{1/2}$; $n_{\mathrm{ext}}\le C(\bar n)$, so $p^{n_{\mathrm{ext}}}\le C_\beta t^{-\beta}$; the \cite[(1.7)]{Ba3} vertices with the finite parts $\dm_{\mathrm{fin}}$, $\dm_{K,k}$ enter through \cref{thm:a06e:T0} and \cref{lem:a06e:T1}. Root the tree at a marked cube ($\le|\widetilde{\mathrm{Col}}|$ choices, including the localised internal vertices of $\dm_k$, $E_k$, \cref{lem:a06f:C21}(iii)) and sum the rest by \cref{lem:a06d:V11}; there are finitely many families of order $\le\bar n$.
\end{proof}

\begin{lemma}[the terms with $W=0$]\label{lem:a06f:C26}
\textup{(a)} The lone \cite[(1.6)]{Ba3} vertex contributes to $P(\Omega_1)-P(\Omega'_1)$ exactly $-\lambda_k\sum_{x\in\mathrm{Col}}\eta^3|\Phi(x)|^4$ (its sum is sharp over $\Omega_1$, \cite[p.~413]{Ba3}). \textup{(b)} The single-line terms contribute $\le C_{\mathrm{SL}}e^2\lambda^{-1/2}p^4t^{1/2}$ per marked cube, hence $\le Ct^{1/2-\beta}|\mathrm{Col}|$ in all.
\end{lemma}

\begin{proof}[Proof of \textup{(b)}]
The order-$1$ objects (both scalar legs external) are $V_8=-e_k\sum_b\eta^3[(D_{\tilde B}\Phi)(b)\cdot q\Phi(b_-)]g_kA'_b$ from \cite[(1.8)]{Ba3}; $P_{12}(y)=a_ke_k\varphi(y)\cdot\sum_{x\in B^k(y)}\eta^3A'(\Gamma^{(k)}_{y,x})qU(\tilde B(\Gamma^{(k)}_{y,x}))\Phi(x)$ from \cite[(1.12)]{Ba3}$\cdot$\cite[(1.14)]{Ba3}$_1$; and $P_{13}(y)$, which is $P_{12}$ with $\varphi(y)$ replaced by $-(Q_k(\tilde B)\Phi)(y)$. The weight $g_k$ forces $y\in\Lambda^{(k-1)'}_7$. By \cite[Lemma~2.4, (2.65)]{Ba2} at every $y\in\Lambda^{(k-1)'}_7$, $\Phi(x)=U(\tilde B(\Gamma^{(k)}_{x,y}))\varphi(y)+\mathrm e(x)$ with $|\mathrm e(x)|\le c_{65}p$; $U(\tilde B(\Gamma_{y,x}))U(\tilde B(\Gamma_{x,y}))=1$ and $\varphi(y)\cdot q\varphi(y)=0$ ($q$ real antisymmetric), so the leading part of $P_{12}$ vanishes and its coefficient is $\le ae_kC_\Phi p\lambda_k^{-1/4}c_{65}p$; the same for $P_{13}$, since $(Q_k(\tilde B)\Phi)(y)=\varphi(y)+O(p)$; for $V_8$, \cite[(2.66)]{Ba2} gives $|D_{\tilde B}\Phi|\le c_{66}p$ on $B^k(\Lambda^{(k-1)'}_7)$, so its coefficient is $\le e_kc_{66}pC_\Phi p\lambda_k^{-1/4}$. The vector line: $A'=\sum_jA'^{(j),\eta}$ with independent Gaussian $A'_j$ of covariance $C^{(j)}$ \cite[p.~414]{Ba3}, and $A'(\Gamma^{(k)}_{y,x})=\sum_jA'^{(j),\eta}(\Gamma^{(j+1)}_{x_{j+1},x})$ \cite[(1.3), p.~412]{Ba3}; its covariance is slice-diagonal, $G^{v}_1+\sum_{j\ge1}S^v_j$ in scale-$k$ units, with $|S^v_j(b,b')|\le Cl_j^{2-d}e^{-\delta_K|b-b'|/l_j}$, $l_j\defeq L^{j-k}$ (\cref{lem:a06e:SL}(b),(d) rescaled; the $j=0$ term with its $L$-dependent constant), so $\sum_{b'}\eta^3|S^v_j(b,b')|\le Cl_j^2$, and the slice-$j$ contour $\Gamma^{(j+1)}$ has $\eta$-length $\le1.25dLl_j$. The three pairings per marked cube: $(V_8,V_8)$: $\le C(e_kp^2\lambda_k^{-1/4})^2\sum_jl_j^2\le Ce^2\lambda^{-1/2}p^4t^{1/2}$; $(V_8,P)$: $\sum_j(1.25dLl_j)Cl_j^2\le C$, the same order; $(P,P')$: per pair of blocks $\sum_j(1.25dLl_j)^2Cl_j^{2-d}e^{-\delta_K\dist/l_j}=\sum_jC'l_je^{-\dots}$ ($d=3$), and the sum over the second block with the decay is $\le C$ (the slice-$j$ field is integrated only over a contour of length $O(l_j)$: no logarithm; the $P$-side bonds are summed through the $\eta^3$-weighted $x\in B^k(y')$, $x\mapsto b'$ being $O(1)$-to-one). Sum the marked cubes over $\widetilde{\mathrm{Col}}$.
\end{proof}

\begin{lemma}[the $c$-number clause]\label{lem:a06f:C27}
\textup{(i)} At zero background on the torus, $c_k(\Omega_1)=-(\Efin+E_{K,k})+a_k(\Omega_1)$ (\cref{lem:a06e:VK0}), so the collar difference within the vertex set of $E_1$ is exactly $0$ at every order $\le\bar n$. \textup{(ii)} At $(\Omega_A,\tilde B)$: $|c_k(\Omega_1)-c_k(\Omega'_1)|\le Ct(1+\ln t^{-1})^5|\mathrm{Col}|\le C_\beta t^{1-\beta}|\mathrm{Col}|$.
\end{lemma}

\begin{proof}
(ii) By \cref{lem:a06f:C22} the difference is a sum over localised vacuum families with a cube meeting $\widetilde{\mathrm{Col}}$ (the vacuum graphs with their $E$-partners, \cite[pp.~420, 422]{Ba3}, and the pair-vertex vacuum graphs). Proposition~1 with no external fields gives $O(1)e_k^{d_v}\lambda_k^{d_s}e^{-\delta_0d}$; a vacuum graph has $d_v\ge2$ or $d_s\ge1$ or contains a \cite[(1.7)]{Ba3} vertex of order $\ge2$ (every expression of $P^{(k)}$ has order $\ge1$ in $e$ unless it is \cite[(1.6)]{Ba3} or \cite[(1.7)]{Ba3}; each $A'$-leg carries one power of $e$; $A'$-legs are contracted in pairs), so $W=2d_v+4d_s\ge4$ and $t^{W/4}\le t$, times $q(t)=(1+\ln t^{-1})^5$ per $\dm_{(4)}$ vertex for the finite counterterm parts (\cref{thm:a06e:T0}, \cref{lem:a06e:T1}); no external field, so no power of $p$. Sum by \cref{lem:a06d:V11} over roots in $\widetilde{\mathrm{Col}}$. At zero background this is \cref{cor:a06d:V13}.
\end{proof}

\begin{proof}[Proof of \cref{thm:a06f:C2}]
\Cref{lem:a06f:C22}(b) splits the difference into the marked families; those with $W\ge2$ are \cref{lem:a06f:C25}, those with $W=0$ are \cref{lem:a06f:C26}, and the field-independent ones are \cref{lem:a06f:C27}. $\eps_0$ enters only through the caps of Step~3(D$_2$) and the ``$e(t)$ sufficiently small'' of \cite[(2.24)]{Ba1} and of Proposition~1; the torus uniformity of Proposition~1's $O(1)$ is the torus side of the propagator estimates of group (A).
\end{proof}

\begin{remark}\label{rem:a06f:HP}
The second sentence of \cite[p.~579]{Ba2} bounds the terms localised in $B^k(\Lambda^{(k)c}_7)$ by $O((L^k\eps)^{\kappa_0})|\Lambda^{(k)}_2\cap\Lambda^{(k)c}_7|$: for $Q^{(k)}$ this is \cref{prop:a06f:H}; for the terms of $P^{(k)}(\Lambda^{(k)}_7,\theta_kA^{(k)},\varphi)$ whose vector leg is gated by $g_k$ there it is \cref{prop:a06g:HP} below, with $\kappa_0=\frac12-\beta$ and the part of the gate inside $\Lambda^{(k)}_7\setminus\Lambda^{(k)}_8$ booked as in \cref{cor:a06g:ledger}; the field-dependent part of the final-step bound \cite[(2.116)]{Ba2} is \cref{prop:a06g:FS}. Both run the mechanism of \cref{thm:a06f:C2}: the gated leg, resp.\ the whole final domain, takes the place of the collar vertex.
\end{remark}

% a06_pgate_a: Step 3(F), fifth part (input from sec/a06_BALtorus_ct8.tex): the new fluctuation field a of Ba2 (2.88)-(2.89)
% at step k (value, derivative and Hoelder bounds) and the geometry of the g_k-gate. Continued in a06_pgate_b (the gated terms
% of P^(k)) and a06_pgate_c (the final step, Ba2 (2.116)).

\subsubsection*{The $g_k$-gated terms of $P^{(k)}$ and the final step \cite[p.~579, (2.116)]{Ba2}}

\emph{Setting.} Step $k$, $1\le k\le K-1$, in scale-$k$ units, with the Setting of \cref{thm:a06f:C2} ($\Omega_A$, $\Omega_1$, $\Omega'_1=B^k(\Lambda^{(k)}_7)$, $\mathrm{Col}=\mathrm{Col}_k$, $\Phi=\phi^{(k)}$, $p=p(t_k)$, $r=r(t_k)$). At step $k+1$ \cite{Ba2} writes $\theta_{k+1}A^{(k)}=\theta_{k+1}A'^{(k)}+\tilde B^{(k)}+\theta_{k+1}B^{(k+1),\eta}$ \cite[(2.88)--(2.89)]{Ba2}, with $A'^{(k)}=a_k\zeta^{(k)}G_kQ_k^*A'$ (\cref{lem:a06d:C62}), $\theta_{k+1}$ the smooth cutoff of $B^k(\Lambda^{(k)}_2)$ changing in a slice of thickness $<M$ \cite[p.~567]{Ba2}, and every term of $\tilde B^{(k)}$ carrying the factor $\theta_{k+1}$. Put
\[
\tilde B\defeq(1-\theta_{k+1})\theta_kA^{(k)}+\theta_{k+1}B^{(k+1),\eta},\qquad a\defeq\theta_{k+1}A'^{(k)}+\tilde B^{(k)},\qquad\text{so that}\qquad\theta_kA^{(k)}=\tilde B+a,
\]
since $\theta_k=1$ on $\supp\theta_{k+1}\subset B^k(\Lambda^{(k-1)'}_7)\subset B^{k-1}(\Lambda^{(k-1)}_2)$ (\cref{lem:a06d:G3}); $\tilde B$ is the background of \cite[p.~582]{Ba2}, regular by \cite[(2.94)--(2.98)]{Ba2}, and $\supp a\subset\supp\theta_{k+1}$, which lies within $M$ of $B^k(\Lambda^{(k)}_2)$, hence in $B^k(\Lambda^{(k)}_1)$ ($\dist(\Lambda^{(k)}_2,\Lambda^{(k)c}_1)\ge r-1>M+1$ under the caps of Step~3(D$_2$)). The gate is $g_k$ ($\supp g_k\subset\Lambda^{(k)}_7$, $g_k=1$ on $\Lambda^{(k)}_8$, $|\partial^\eta g_k|\le c_g$, as in \cref{prop:a06f:H}). The cutoffs are taken, as the printed ``smooth'' allows, from fixed classes: $\theta_{k+1}$ a fixed smooth function of the mollified distance to $B^k(\Lambda^{(k)}_2)$, so that $|\partial^\eta\theta_{k+1}|\le c_\theta$ and $|\partial^\eta\partial^\eta\theta_{k+1}|\le c_{\theta2}$ uniformly in $k$ and $\eta$ (a smooth function of the lattice sup-distance would have second differences of order $\eta^{-1}$ along the edges of the slice); and $\zeta^{(k)}$ with $|\zeta^{(k)}|\le1$, $|\partial^\eta_x\zeta^{(k)}|\le1$, $|\partial^\eta_x\partial^\eta_x\zeta^{(k)}|\le2$ (the cut-off of \cref{thm:a06:BALA}(b) is so chosen: a fixed smooth profile of the mollified distance to $y$, as for $\theta_{k+1}$, over its transition region of width $\frac15r-2M$, which is $\ge2$ for $r\ge10M+10$, a cap on $\eps_0$ free of $k$). Put $\delta'\defeq\min(\delta_*,\delta_0)$, $S'_*\defeq\sum_{n\in\Z^3}e^{-\delta'|n|/2}$, and $C'(\delta')$ for the constant $C'$ of \cref{lem:a06d:C62} with $\delta'$ in place of $\delta_*$.

\begin{lemma}[the field $a$]\label{lem:a06g:A}
On $\supp\theta_{k+1}$, for $t_k\le t_0$, where $t_0\in(0,1]$ is fixed by $C_\kappa t_0^\kappa\le C_{2.91}b_0$, $C'_\kappa t_0^\kappa\le b_0$ and $C^H_\kappa t_0^\kappa\le b_0$ ($C_\kappa$, $C'_\kappa$ the constants of (R2) of Step~3(C$_2$), $C^H_\kappa$ that of (c) below; a cap on $\eps_0$, free of $k$; $p(t)\ge b_0$): \textup{(a)} $|a|\le C_{2.91}p+C_\kappa t_k^\kappa\le2C_{2.91}p$; \textup{(b)} for every fine bond $b=(x,\mu)$ with $x$ or $x+\eta e_\mu$ in $\supp\theta_{k+1}$, $|\partial^\eta_\mu a(x)|\le C'_{2.91}p+C'_\kappa t_k^\kappa\le(C'_{2.91}+1)p$, with
\[
C'_{2.91}\defeq c_\theta C_{2.91}+a(C_1e^{\delta_*}+c_0e^{\delta_0})S'_*(C_{2.86}+C'(\delta'));
\]
\textup{(c)} for parallel bonds $b,b'$ within distance $M$ of $\supp\theta_{k+1}$ with $|b-b'|\le1$, $|\partial^\eta a(b)-\partial^\eta a(b')|\le C^H_{2.91}p\,|b-b'|^{\alpha_0}$. The constants are free of $k$, $\eps$, $\eps_0$, the torus and the history. So $a$ satisfies the hypotheses on an external vector field of \cite[p.~412]{Ba3} with $\norm{\cdot}_{1,\alpha_0}$-norm $\le C_ap$, $C_a\defeq2C_{2.91}+C'_{2.91}+1+C^H_{2.91}$.
\end{lemma}

\begin{proof}
(a) is \cref{lem:a06d:C62} for $\theta_{k+1}A'^{(k)}$ and (R2) of Step~3(C$_2$) for $\tilde B^{(k)}$ \cite[(2.90)]{Ba2}; $C_\kappa t^\kappa\le C_{2.91}p$ for $t\le t_0$.

(b) \emph{Lattice distances.} Fine points are at mutual distances in $\eta\Z$, and the slice has thickness $<M$, so $\supp\theta_{k+1}$ lies within $M-\eta$ of $B^k(\Lambda^{(k)}_2)$ and both $x$ and $x+\eta e_\mu$ lie within $M$ of it; \cref{lem:a06d:C62} and its proof apply at both points (the proof bounds $A'^{(k)}$ itself there, without the factor $\theta_{k+1}$). \emph{Leibniz rule.} $\partial_\mu(\theta_{k+1}A'^{(k)})(x)=(\partial_\mu\theta_{k+1})(x)A'^{(k)}(x+\eta e_\mu)+\theta_{k+1}(x)(\partial_\mu A'^{(k)})(x)$. The first term is $\le c_\theta C_{2.91}p$. In the second,
\[
\partial_\mu A'^{(k)}(x)=a_k\sum_y\big[(\partial_{\mu,x}\zeta^{(k)})(x,y)(G_kQ_k^*)(x+\eta e_\mu,y)+\zeta^{(k)}(x,y)(\partial_\mu G_kQ_k^*)(x,y)\big]A'(y);
\]
both kernels are $\le(C_1e^{\delta_*}+c_0e^{\delta_0})e^{-\delta'|x-y|}$, by $|\partial\zeta|\le1$, \cref{thm:a06:A5} ($|\cdot|_1\ge|\cdot|$ and one bond shift) and the derivative clause of \cite[Prop.~2.2, (2.58)]{Ba2} (\cite[Thm., (1.9)--(1.10)]{Ba4}, at zero background on the torus, where the depth restriction is void, \cref{lem:a06A:torus}); both vanish unless $|x-y|<r-2M+\eta$, and such $y$ lie in $\Lambda^{(k)}_0$ by the layer count of \cref{lem:a06d:C62} ($2(r-c_b)>r-2M+M+1$). Split the $y$-sum as there: on $\Lambda^{(k)}_1$, \cref{lem:a06d:C61} gives $\le(\dots)S'_*C_{2.86}p$; on $\Lambda^{(k)}_0\setminus\Lambda^{(k)}_1$, $|x-y|\ge r-c_b-M-1$ and the factor $e^{-\delta'|x-y|}$ with $E(1,\delta'/2)$ turn $p/(\lmu t)$ into $C'(\delta')p$. With $|\theta_{k+1}|\le1$ and $a_k\le a$ the second term is $\le a(C_1e^{\delta_*}+c_0e^{\delta_0})S'_*(C_{2.86}+C'(\delta'))p$. For $\tilde B^{(k)}$ the derivative bound is (R2).

(c) Write $\partial(\theta_{k+1}A'^{(k)})=(\partial\theta_{k+1})\,\mathsf SA'^{(k)}+\theta_{k+1}\partial A'^{(k)}$ ($\mathsf S$ the bond shift) and telescope the difference at $b,b'$ term by term. The factors $\partial\theta_{k+1}$, $\theta_{k+1}$ and the kernel factor $\partial\zeta^{(k)}$ are Lipschitz on the $\eta$-lattice with constants $c_{\theta2}$, $c_\theta$ and $2$; at separation $|b-b'|\ge\eta$ a Lipschitz bound is $\le$ the constant times $|b-b'|^{\alpha_0}$, since $|b-b'|\le1$; their cofactors ($A'^{(k)}$, $\partial A'^{(k)}$ and the kernel sums) are bounded by \cref{lem:a06d:C62} and by (b), the kernel $(G_kQ_k^*)(\cdot+\eta e_\mu,y)$ of the $\partial\zeta^{(k)}$ term being itself Lipschitz in $x$ with constant $c_0e^{\delta_0}e^{-\delta_0\dist(\cdot,y)}$ by the derivative clause of \cite[(2.58)]{Ba2} (a difference along a lattice path of length $|b-b'|\le1$), summed over $y$ as in (b). The $\theta$-factors are telescoped first, so each cofactor is either evaluated at a bond of $\supp\theta_{k+1}$ or is a difference of two kernel sums of the type of (b) at points within $M+1$ of $B^k(\Lambda^{(k)}_2)$; the bounds of (b) hold there, since the kernel bounds are global on the torus and the layer count with $M+1$ in place of $M$ keeps the $y$-range in $\Lambda^{(k)}_0$. The term $\zeta^{(k)}(\partial G_kQ_k^*)A'$: its variation in $b$ is the H\"older clause of \cite[Prop.~2.1, (2.24)]{Ba1} (\cite[Thm., (1.10)]{Ba4}) for $\partial G_kQ_k^*\delta_y$, $|(\partial G_kQ_k^*)(b,y)-(\partial G_kQ_k^*)(b',y)|\le c_0|b-b'|^{\alpha_0}e^{-\delta_0\min(\dist(b,y),\dist(b',y))}$; sum over $y$ as in (b). For $\tilde B^{(k)}$: by \cite[(2.89)]{Ba2} each of its three terms is $\theta_{k+1}$ times a kernel sum of the same type ($\zeta^{(k)}G_kQ_k^*$, its tail $(1-\zeta^{(k)})G_kQ_k^*$, or $\zeta^{(k)}G_kQ_k^*$ composed with the covariance difference of \cref{lem:a06A:CLdiff}) applied to a block field bounded by $Cp/(\lmu t)$ on $\Lambda^{(k)'}_0$ (the restrictions \cite[(2.17), (2.55)]{Ba2}); the same computation bounds its H\"older quotient, and the separation $\ge r-2c_b-M$ between the evaluation points in $\supp\theta_{k+1}$ and the support of the field, through $E(1,\delta/2)$, gives $C^H_\kappa t_k^\kappa$ in place of $p$, as in (R2).
\end{proof}

\begin{lemma}[the gate]\label{lem:a06g:BAD}
Let $\mathrm{Bad}_k$ be the set of scale-$k$ unit cubes meeting $\supp((1-g_k)a)$. Then $\mathrm{Bad}_k\subset N_M(B^k(\Lambda^{(k)}_2))\setminus B^k(\Lambda^{(k)}_8)\subset\Omega_1$, every cube of $\mathrm{Bad}_k$ lies within $M$ of a cube of $\Lambda^{(k)}_2\cap\Lambda^{(k)c}_8$, and
\[
|\mathrm{Bad}_k|\le(2M+1)^3|\Lambda^{(k)}_2\cap\Lambda^{(k)c}_8|\le(2M+1)^3\big(|\Lambda^{(k)}_2\cap\Lambda^{(k)c}_7|+|\Lambda^{(k)}_7\setminus\Lambda^{(k)}_8|\big),\qquad\Lambda^{(k)}_2\cap\Lambda^{(k)c}_7\subset\mathrm{Col}_k.
\]
\end{lemma}

\begin{proof}
$(1-g_k)a$ vanishes outside $\supp\theta_{k+1}\subset N_M(B^k(\Lambda^{(k)}_2))$ and on $B^k(\Lambda^{(k)}_8)$, where $g_k=1$. Let $y\in\mathrm{Bad}_k$ and let $x$ be a cube of $B^k(\Lambda^{(k)}_2)$ within $M$ of $y$. If $y\in\Lambda^{(k)}_2$, then $y\in\Lambda^{(k)}_2\cap\Lambda^{(k)c}_8$ itself. If $y\notin\Lambda^{(k)}_2$, then $\dist(x,\Lambda^{(k)c}_2)\le M+1<6(r-1)\le\dist(\Lambda^{(k)}_8,\Lambda^{(k)c}_2)$, so $x\notin\Lambda^{(k)}_8$ and $x\in\Lambda^{(k)}_2\cap\Lambda^{(k)c}_8$. A cube of $\Lambda^{(k)}_2\cap\Lambda^{(k)c}_8$ has at most $(2M+1)^3$ cubes within $M$ of it. The inclusions: $\Lambda^{(k)}_8\subset\Lambda^{(k)}_7$, and $\Lambda^{(k)}_2\subset\Lambda^{(k)}_0\subset\Lambda^{(k-1)'}_7$ \cite[pp.~566, 570]{Ba2}, so $\Lambda^{(k)}_2\cap\Lambda^{(k)c}_7\subset\Lambda^{(k-1)'}_7\cap\Lambda^{(k)c}_7=\mathrm{Col}_k$ and $N_M(B^k(\Lambda^{(k)}_2))\subset B^k(\Lambda^{(k)}_1)\subset\Omega_1$.
\end{proof}

% a06_pgate_b: Step 3(F), fifth part, continued (input from sec/a06_BALtorus_ct8.tex): the terms of P^(k) whose vector leg is
% gated by g_k outside Lambda_8^(k) (the P^(k)-half of the second sentence of Ba2 p. 579), and their place in the ledger.

In the Setting above, expand the step-$k$ polynomial at the background $\theta_kA^{(k)}=\tilde B+a$ in powers of the new field $a$ \cite[(3.15), (3.44)]{Ba1}, \cite[p.~579]{Ba2}:
\[
P^{(k)}(\Omega'_1,\tilde B+a,\varphi)=\sum_{n_a=0}^{\bar n}P^{(k)}_{n_a}[a^{n_a}]+R^{(k)}_a,
\]
where $P^{(k)}_{n_a}[a^{n_a}]$ is the sum of the localised renormalised families of $P^{(k)}(\Omega'_1,\tilde B,\varphi)$ carrying $n_a$ external vector legs $a$: the vertices \cite[(1.8)--(1.11)]{Ba3} with $n'\ge1$, and the $a$-insertions on the scalar lines $G_k(\Omega_A,\tilde B+a)$ and on the external fields $\Phi(\tilde B+a)$ of \cite[(2.56)]{Ba2}; $R^{(k)}_a$ collects the $R$-vertices and the orders $>\bar n$ in $a$. On p.~579 the new legs are localised by $g_k$. Splitting every leg as $a=g_ka+(1-g_k)a$ (multilinearity), put
\[
T^{(k)}_{\mathrm{gate}}\defeq\sum_{n_a=1}^{\bar n}\big(P^{(k)}_{n_a}[a^{n_a}]-P^{(k)}_{n_a}[(g_ka)^{n_a}]\big),
\]
the sum of the terms with at least one leg carrying $(1-g_k)a$, i.e.\ with a vector leg localised where $g_k\ne1$: in $B^k(\Lambda^{(k)c}_7)$ or in $B^k(\Lambda^{(k)}_7\setminus\Lambda^{(k)}_8)$. If $P^{(k)}(\Lambda^{(k)}_7,\theta_kA^{(k)},\varphi)$ in \cite[(2.57)]{Ba2} is understood as the polynomial already built with the gated field $g_ka$, there is no such term; otherwise the gating is a move, whose cost is the following proposition.

\begin{proposition}[{the terms of $P^{(k)}$ removed by $g_k$; the $P^{(k)}$-half of the second sentence of \cite[p.~579]{Ba2}}]\label{prop:a06g:HP}
For $1\le k\le K-1$, under \cite[(2.55)]{Ba2} and the caps of Step~3(D$_2$), for every $\beta>0$,
\[
|T^{(k)}_{\mathrm{gate}}|\le C_{\mathrm{HP}}t_k^{1/2-\beta}|\mathrm{Bad}_k|\le C_{\mathrm{HP}}(2M+1)^3t_k^{1/2-\beta}\big(|\mathrm{Col}_k|+|\Lambda^{(k)}_7\setminus\Lambda^{(k)}_8|\big),
\]
with $C_{\mathrm{HP}}$ depending on the data of \cref{thm:a06f:C2} and on $c_\theta,c_{\theta2},c_g$, and free of $\eps$, $k$, $\eps_0$, the torus and the history. The remainder $R^{(k)}_a$ (the $R$-vertices and the orders $>\bar n$ in $a$) is $O(t_k^\kappa)|\Lambda^{(k)}_2|$ for every $\kappa$ (\cref{prop:a06e:B12}(iii)), the first half of the sentence.
\end{proposition}

\begin{proof}
\emph{Step 1: the weight.} Count every $a$-leg in $d_v$: the $a$-dependence of each parallel transporter is $U(\tilde B+a)=U(\tilde B)U(a)$ (the group is abelian), so each $a$-insertion carries its own factor $e_k$, and $W\defeq2d_v+4d_s-n_{\mathrm{ext}}$ is defined for a family with $a$-legs as in \cref{lem:a06f:C24}, the $a$-legs included in $d_v$. Inserting one $a$-leg into a family raises $W$ by at least $2$: on an existing gauge vertex, $d_v(v)$ rises by $1$ and $X(v)$ by $2$; as a new vertex of the type \cite[(1.8)]{Ba3} with $n=1$ (one vector leg, two scalar legs) splitting a scalar line, whether an internal line or the line $G_k(\Omega_A,\tilde B+a)$ of an external field $\Phi$, $X(\text{new})=2-2=0$ and $l_s$ rises by $1$, while $n_{\mathrm{ext}}$ is unchanged (the part of the split line beyond the new vertex is the external field $\Phi(\tilde B)$); as an insertion on the block average $Q_k^*(\tilde B+a)$ of an external field, $d_v$ rises by $1$ and nothing else changes: the external field becomes $a_kG_k(\Omega_A,\tilde B)f_a$ with the source $f_a\defeq[Q_k^*(\tilde B+a)-Q_k^*(\tilde B)]\Lambda^{(k-1)'}_6\varphi$, $\norm{f_a}_\infty\le e_k(L-1)d\sup|a|\sup|\varphi|$ (the transporter of $Q_k^*$ runs along a contour of $\eta$-length at most $(L-1)d$ in scale-$k$ units, \cite[p.~559]{Ba2}, and $q$ is antisymmetric), which the proof of \cref{lem:a06f:C23} bounds by $e_k(L-1)d\,2C_{2.91}p\cdot C_\Phi p\lambda_k^{-1/4}$, the factor $e_kC_ap$ of Step~3 up to $O(1)$. By \cref{lem:a06f:C24} every family of $P^{(k)}$ has $W\ge0$, so every term of $T^{(k)}_{\mathrm{gate}}$ has $W\ge2$ ($W$ stays even). In particular the $W=0$ objects of \cref{lem:a06f:C26} (the lone quartic and the single-line terms) occur in $T^{(k)}_{\mathrm{gate}}$ only with an $a$-leg attached, hence with $W\ge2$, and the cancellation of \cref{lem:a06f:C26}(b) is not needed.

\emph{Step 2: the fields.} The external vector field $a$ obeys \cref{lem:a06g:A}: $|a|\le2C_{2.91}p$ on $\supp\theta_{k+1}$ \cite[(2.90)--(2.91)]{Ba2}, $|\partial^\eta a|\le(C'_{2.91}+1)p$, H\"older quotient $\le C^H_{2.91}p$; so it is an admissible external vector field of \cite[p.~412]{Ba3} with norm $\le C_ap$. The external scalar fields are those of \cref{lem:a06f:C23}. The background $\tilde B$ is regular by \cite[(2.94)--(2.98)]{Ba2}, as in \cref{prop:a06e:B12}(iii).

\emph{Step 3: the bound per family.} Proposition~1 with the external vector field $a$ (the legs of \cite[(1.10)]{Ba3} with $n'\ge1$; its $O(1)$ is independent of $\eps$, $k$, the domains and the background, \cite[p.~421]{Ba3}) bounds each localised family by
\[
O(1)\,e^{d_v}\lambda^{d_s-n_{\mathrm{ext}}/4}(C_\Phi p)^{n_{\mathrm{ext}}}(C_ap)^{n_a}\,t^{W/4}e^{-\delta_0d(\{\square\})},
\]
by the computation of \cref{lem:a06f:C25} in $d=3$ ($e_k^{d_v}\lambda_k^{d_s}(p\lambda_k^{-1/4})^{n_{\mathrm{ext}}}=e^{d_v}\lambda^{d_s-n_{\mathrm{ext}}/4}p^{n_{\mathrm{ext}}}t^{W/4}$, the $e_k$ of the $a$-legs included in $e^{d_v}t^{d_v/2}$); the \cite[(1.7)]{Ba3} vertices with the finite parts $\dm_{\mathrm{fin}}$, $\dm_{K,k}$ enter through \cref{thm:a06e:T0} and \cref{lem:a06e:T1}. By Step~1, $t^{W/4}\le t^{1/2}$; $n_{\mathrm{ext}}+n_a\le C(\bar n)$, so $p^{n_{\mathrm{ext}}+n_a}\le C_\beta t^{-\beta}$ ($P(\beta,n)<\infty$); hence each family is $\le C_\beta t^{1/2-\beta}e^{-\delta_0d(\{\square\})}$.

\emph{Step 4: the sum.} Root the tree of each family at the cube of a leg carrying $(1-g_k)a$: at most $|\mathrm{Bad}_k|$ choices; sum the remaining cubes by \cref{lem:a06d:V11}; there are finitely many families of order $\le\bar n$ with $n_a\le\bar n$. This gives the first bound, and \cref{lem:a06g:BAD} the second.

\emph{Step 5: what is not gated.} The $R$-vertices and the orders $>\bar n$ in $a$ are not terms of $T^{(k)}_{\mathrm{gate}}$; they are bounded in \cref{prop:a06e:B12}(iii) by $O(t_k^\kappa)|\Lambda^{(k)}_2|$, Proposition~1 with the field restrictions and $\bar n+1$ legs $a$.
\end{proof}

\begin{corollary}[the two parts in the ledger]\label{cor:a06g:ledger}
The part of $T^{(k)}_{\mathrm{gate}}$ with a gated leg in $B^k(\Lambda^{(k)c}_7)$ is $\le C_{\mathrm{HP}}(2M+1)^3t_k^{1/2-\beta}|\mathrm{Col}_k|$: a collar term of the form of \cite[(2.57)]{Ba2}'s remainder, which enters \cite[(2.117)]{Ba2} with the collar terms (it is the $P^{(k)}$-part of the printed $O((L^k\eps)^{\kappa_0})|\Lambda^{(k)}_2\cap\Lambda^{(k)c}_7|$). The part with a gated leg in $B^k(\Lambda^{(k)}_7\setminus\Lambda^{(k)}_8)$ is $\le C_{\mathrm{HP}}(2M+1)^3|\Lambda^{(k)c}_8|$ ($t_k^{1/2-\beta}\le1$): a term of the form of \cref{prop:a06f:D2C}'s moves and of \cref{prop:a06f:H}'s bound, which enters the augmented exponent of \cref{thm:a06f:R7} with the same sign and in the same place. \Cref{thm:a06f:R7} holds for every $k$-free $c_8\ge0$; the three terms of this kind are carried together by taking
\[
c_8\ \ge\ Q_1(C_{\mathrm{D22}}+C_{\mathrm{D23}})+(2M+1)^3\big(\sup_{0<t\le1}F_H(t)+C_{\mathrm{HP}}\big),
\]
a constant free of $k$, $\eps$, $K$, the torus and the history, fixed before $\eps_0$, so that (E.ent) of \cref{thm:a06f:R7} is one cap on $\eps_0$ as before.
\end{corollary}

\begin{proof}
\Cref{lem:a06g:BAD} splits $\mathrm{Bad}_k$ into the cubes within $M$ of $\Lambda^{(k)}_2\cap\Lambda^{(k)c}_7\subset\mathrm{Col}_k$ and those within $M$ of $\Lambda^{(k)}_7\setminus\Lambda^{(k)}_8\subset\Lambda^{(k)c}_8$; the rooting of Step~4 is done in the respective set. The bookkeeping is that of the proof of \cref{prop:a06f:D2C}, and \cref{lem:a06f:C141} carries the terms through \cite[Sect.~3A]{Ba2}; $F_H$ is bounded on $(0,1]$ (\cref{prop:a06f:H}).
\end{proof}

% a06_pgate_c: Step 3(F), fifth part, end (input from sec/a06_BALtorus_ct8.tex): the final step, the field-dependent part of
% Ba2 (2.116), as Theorem thm:a06f:C2 at k = K with an empty inner domain.

\emph{The final step.} Let $K$ be the last step, $l_K=t_K=L^K\eps\in(\eps_0/L,\eps_0]$ (\cref{thm:a06:F}), in scale-$K$ units, with the Setting of \cref{thm:a06f:C2} at $k=K$: $\Omega_A=B^K(\Lambda^{(K-1)'}_2)$, $\Omega_1=B^K(\Lambda^{(K-1)'}_7)$, $\tilde B=A^{(K)}=\theta_KA^{(K)}$ on $\Omega_A$, $\Phi=\phi^{(K)}$ of \cite[(2.56)]{Ba2} at $k=K$, under \cite[(2.55)]{Ba2} at step $K$. No set $\Lambda^{(K)}_i$ is defined at the final step: \cite[(2.116)]{Ba2} bounds $P^{(K)}$ on its whole domain $\Omega_1$, and its right side is the volume $|\Lambda^{(K-1)'}_7|$ of that domain, as in the last term of \cite[(2.117)]{Ba2}. The chain is run with $E_0+E_K$ (\cref{lem:a06:CF}), so no step polynomial contains $\Efin$, which is the global factor bounded by \cref{thm:a06:F}.

\begin{proposition}[{the final step, \cite[(2.116)]{Ba2}}]\label{prop:a06g:FS}
For every $\beta>0$,
\[
P^{(K)}(\Omega_1,\tilde B,\varphi)=-\lambda_K\sum_{x\in\Omega_1}\eta^3|\Phi(x)|^4+R_{\mathrm{field}}+c_K(\Omega_1),\qquad|R_{\mathrm{field}}|\le C_\beta t_K^{1/2-\beta}|\Lambda^{(K-1)'}_7|,
\]
with $c_K(\Omega_1)$ the field-independent part, $|c_K(\Omega_1)|\le Ct_K(1+\ln t_K^{-1})^5|\Lambda^{(K-1)'}_7|\le C_\beta t_K^{1-\beta}|\Lambda^{(K-1)'}_7|$. Hence $P^{(K)}(\Omega_1,\tilde B,\varphi)\le C_\beta t_K^{1/2-\beta}|\Lambda^{(K-1)'}_7|$, which is \cite[(2.116)]{Ba2} with $\kappa_0=\frac12-\beta$. The constants are those of \cref{thm:a06f:C2}, free of $\eps$, $K$, $\eps_0$, the torus and the history.
\end{proposition}

\begin{proof}
Run the proof of \cref{thm:a06f:C2} at $k=K$ with the inner domain $\Omega'_1\defeq\emptyset$.

\emph{Step 1: the empty domain.} By \cref{lem:a06f:C21} the vertex domain enters only through $\mathbf 1_{\Omega'_1}$ on the vertices \cite[(1.6)--(1.7)]{Ba3}, through the gate on the $A'$-legs ($\supp g\subset\Omega'_1$) and through the localised counterterms. With $\Omega'_1=\emptyset$ every vertex vanishes and every $A'$-leg with it, so $P^{(K)}(\emptyset,\tilde B,\varphi)$ is field-independent and equals $c_K(\emptyset)$, which is $0$: on the chain run with $E_0+E_K$, $c_k(\Omega')=-E_{K,k}+a_k(\Omega')$ (\cref{lem:a06e:VK0}), $E_{K,K}=E_K-E_K=0$, and $a_K(\emptyset)$ has no pair vertex to place. In \cref{lem:a06f:C22}, $D_{\mathrm{good}}=\emptyset$ and $\widetilde{\mathrm{Col}}=\Omega_1$: every localised family is marked, and $P^{(K)}(\Omega_1)=P^{(K)}(\Omega_1)-P^{(K)}(\emptyset)$ is the sum over all localised families, rooted at a cube of $\Omega_1$ ($|\Omega_1|=|\Lambda^{(K-1)'}_7|$ choices).

\emph{Step 2: the three kinds of family.} The families with $W\ge2$ are bounded by \cref{lem:a06f:C25} at $k=K$: Proposition~1 on the lines $G_K(\Omega_A,\tilde B)$ with the external scalar fields of \cref{lem:a06f:C23} at $k=K$ (which uses \cite[(2.55)]{Ba2} at step $K$ and \cite[(2.24)--(2.25)]{Ba1} at points of $\Omega_1$, at depth $\ge r(t_K)-2\ge R_0$ in $\Omega_A$, $\tilde B$ regular by \cite[(2.60)]{Ba2}), the \cite[(1.7)]{Ba3} vertices through \cref{thm:a06e:T0} and \cref{lem:a06e:T1}, the sum by \cref{lem:a06d:V11} with the root in $\Omega_1$: in all $\le C_\beta t_K^{1/2-\beta}|\Lambda^{(K-1)'}_7|$. The families with $W=0$ are \cref{lem:a06f:C26}: the lone \cite[(1.6)]{Ba3} vertex, whose sum over the vertex domain is sharp, gives exactly $-\lambda_K\sum_{x\in\Omega_1}\eta^3|\Phi(x)|^4$; the single-line terms give $\le C_{\mathrm{SL}}e^2\lambda^{-1/2}p^4t_K^{1/2}$ per cube of $\Omega_1$ by \cite[Lemma~2.4, (2.65)--(2.66)]{Ba2} at every $y\in\Lambda^{(K-1)'}_7$ (\cref{thm:a06:BALA}(c)) and the slice kernels of \cref{lem:a06e:SL}, in all $\le Ct_K^{1/2-\beta}|\Lambda^{(K-1)'}_7|$. The field-independent families are \cref{lem:a06f:C27}(ii) with $\Omega'_1=\emptyset$: the vacuum families with a cube meeting $\Omega_1$ have $W\ge4$, and $|c_K(\Omega_1)|=|c_K(\Omega_1)-c_K(\emptyset)|\le Ct_K(1+\ln t_K^{-1})^5|\Lambda^{(K-1)'}_7|$; at zero background this is \cref{cor:a06d:V13}, and on the chain of \cref{lem:a06:CF} it is the bound of \cref{thm:a06e:VK}(a) on $a_K(\Omega_1)$, whose proof does not use $k<K$.

\emph{Step 3: the inputs at $k=K$.} Every input of \cref{thm:a06f:C2} is available at the final step: \cite[(2.55)]{Ba2} holds at every step, $K$ included; Proposition~1's $O(1)$ does not depend on the step; the layer distances used are those of the step-$(K-1)$ sets, $\dist(\Omega_1,\Omega_A^c)\ge r(t_K)-2$; and no set $\Lambda^{(K)}_i$ enters, since the inner domain is empty. The one-sided bound follows by dropping the quartic term, which is $\le0$.
\end{proof}

The same bound, with $|\Lambda^{(K-1)'}_7|$ replaced by $L^3|T^{(K)}|\ge|\Lambda^{(K-1)'}_7|$, is the $k=K$ term of the exponents of \cite[(3.42)]{Ba2} as \cref{prop:a06f:RS} carries it. This completes the field-dependent part of the second sentence of \cite[p.~579]{Ba2} (\cref{prop:a06g:HP}) and of the final step (\cref{prop:a06g:FS}); with \cref{thm:a06f:C2} (the collar), the three transfers of the mechanism of \cite[Prop.~2.1]{Ba2} that \cite[p.~570]{Ba2} announces (``similar conclusions hold for other expressions which will be included in the action in the later stages of the procedure'') are written.

